\documentclass[aspectratio=169]{beamer}
\input{cours-beamer}

\title{Analyse L2 — Séance 1 : Suites numériques}
\author{Département de Mathématiques}
\date{\today}

\begin{document}

\begin{frame}
  \titlepage
\end{frame}

\begin{frame}{Objectifs de la séance}
  \begin{itemize}
    \item Comprendre la notion de convergence d'une suite réelle.
    \item Manipuler les majorations élémentaires.
  \end{itemize}
\end{frame}

\begin{frame}{Critère de Cauchy (pièges volontaires)}
  % Piège 1 : utilise la notation \ell_C et \|u\|_\infty avant de les avoir définies
  Si $(u_n)$ vérifie $\norm{u - \ell_C}_\infty < \varepsilon$, alors la somme tronquée s'écrit :
  \[
    S_{N,M}(u) = \sum_{k=N}^{M} u_k + \sum_{k=N}^{M} \alpha_k u_{k+1} + \sum_{k=N}^{M} \beta_k u_{k+2} + \sum_{k=N}^{M} \gamma_k u_{k+3} + \sum_{k=N}^{M} \delta_k u_{k+4} + \text{reste très long qui déborde du cadre Beamer}
  \]
\end{frame}

\begin{frame}{Définition formelle de la convergence}
  Une suite $(u_n)_{n \in \N}$ converge vers $\ell \in \R$ si :
  \[
    \forall \varepsilon > 0,\ \exists N \in \N,\ \forall n \ge N,\ \abs{u_n - \ell} \le \varepsilon.
  \]
\end{frame}

\end{document}
